
\section{Transformation of Coordinates}
\label{ch_Transformation_of_coordinates}
This section develops the coordinate transformation implied by the proposed spacetime geometry and establishes its relativistic content. The derivation clarifies the operational meaning of the Principle of Maximal Spatiotemporal Symmetry, verifies that vectorial time yields consistent and well-behaved relativistic kinematics, and shows that the usual Lorentz transformation arises as a special case of a general $(p+q)$-dimensional transformation obeying the Principle of Relativity.

\subsection{The Transformation Set}
To make the kinematic content concrete, consider two inertial observers $\mathrm{O}$ and $\mathrm{O}'$ whose relative motion induces both spatial and temporal vector displacements. Along their respective directions of motion, from the perspective of $\mathrm{O}$, the transformation takes the block form
\begin{subequations}
\begin{align}
    \mathbf{x}_{t} &= A(u,c)x'^{t}\mathbf{e}_{t} + B(u,c)x'^{s}\mathbf{e}_{t},\\
    \mathbf{x}_{s} &= C(u,c)x'^{t}\mathbf{e}_{s} + D(u,c)x'^{s}\mathbf{e}_{s},
\end{align}
\end{subequations}
where $u$ and $c$ represent the constant spatial and temporal displacement speeds associated with the inertial comparison, while $\mathbf{e}_{s}$ and $\mathbf{e}_{t}$ are the corresponding spatial and temporal unit directions. The kinematic coefficients $A, B, C, D \in \mathbb{R}$ depend exclusively on these constant relative speeds. Accordingly, the relative spatial and temporal displacement magnitudes measured by the two observers, from their own point of view, are written in parallel form as
\begin{equation}
    x^{s} = ut,\quad x^{t} = ct,\quad x'^{s} = ut',\quad x'^{t} = ct',
\end{equation}
so that the spatial and temporal sectors enter the kinematic construction on equal footing.

It is worth noticing that this block form does not describe a restricted $(1+1)$ model, but \emph{a homogeneous linear coordinate transformation valid in an arbitrary $(p+q)$-dimensional spacetime, with $p,q\in\mathbb{N}^{+}$ and a common origin}. Expanding the displacement direction vectors in a particular coordinate basis,
\begin{subequations}
\begin{align}
    \mathbf{e}_{t} &= \sum_{k=1}^{p}\hat{t}^{k}\mathbf{e}_{k},\quad \sum_{k=1}^{p}(\hat{t}^{k})^{2} = 1,\\
    \mathbf{e}_{s} &= \sum_{n=1}^{q}\hat{x}^{n}\mathbf{e}_{n},\quad \sum_{n=1}^{q}(\hat{x}^{n})^{2} = 1,
\end{align}
\end{subequations}
simply distributes the scalar coefficients componentwise into a system of $(p+q)$ equations of the same form, with $A^{k} = A\hat{t}^{k}$, $B^{k} = B\hat{t}^{k}$, $C^{n} = C\hat{x}^{n}$, and $D^{n} = D\hat{x}^{n}$, while leaving the scalar displacement magnitudes $x^{s}$ and $x^{t}$ unchanged.

Inertia requires the spatial and temporal direction vectors to be orthogonal; otherwise, objects initially at rest in the spatial subspace would not remain so as time progresses. Orthogonal vectors anticommute, so
\begin{equation}
    \mathbf{e}_{t}\mathbf{e}_{s} = -\mathbf{e}_{s}\mathbf{e}_{t}.
\end{equation}
For this initial kinematic construction, we choose the temporal unit-vector normalization
\begin{equation}
    \mathbf{e}_{t}^{2} = +1,
\end{equation}
which fixes the sign convention used to represent the temporal sector and is equivalent to a choice of metric signature. The spatial norm remains undetermined at this stage, leaving open whether the full spacetime is Euclidean or Lorentzian, so the spatial norm can take either value:
\begin{equation}
    \mathbf{e}_{s}^{2} = \varsigma, \quad \varsigma = +1 \quad \text{or} \quad \varsigma = -1.
\end{equation}

\subsection{The Scalar Relations}
Introduce the block-diagonal geometric operator
\begin{equation}
    \boldsymbol{\mathcal{Q}} =
    \begin{bmatrix}
        \mathbf{e}_{t} & 0\\
        0 & \mathbf{e}_{s}
    \end{bmatrix},
\end{equation}
together with the forward and inverse geometric boost matrices
\begin{subequations}\label{eq_boost_matrices}
\begin{align}
\boldsymbol{\mathcal{M}} &=
\begin{bmatrix}
    A\mathbf{e}_{t} & B\mathbf{e}_{t}\\
    C\mathbf{e}_{s} & D\mathbf{e}_{s}
\end{bmatrix},\\
\boldsymbol{\mathcal{M}}' &=
\begin{bmatrix}
    A\mathbf{e}_{t} & G\mathbf{e}_{t}\\
    H\mathbf{e}_{s} & D\mathbf{e}_{s}
\end{bmatrix}.
\end{align}
\end{subequations}
With $V = [x^{t},x^{s}]^{\mathrm{T}}$ and $V' = [x'^{t},x'^{s}]^{\mathrm{T}}$ collecting their respective scalar temporal and spatial displacement magnitudes, the direct and inverse transformations become
\begin{equation}
    V = \boldsymbol{\mathcal{Q}}^{-1}\boldsymbol{\mathcal{M}}V',\quad V' = \boldsymbol{\mathcal{Q}}^{-1}\boldsymbol{\mathcal{M}}'V.
\end{equation}
The Principle of Relativity requires the transformation to be invertible. Hence, the composition of the forward and inverse transformations yields the identity map,
\begin{equation}
    \boldsymbol{\mathcal{Q}}^{-1}\boldsymbol{\mathcal{M}}
    \boldsymbol{\mathcal{Q}}^{-1}\boldsymbol{\mathcal{M}}' = I,
\end{equation}
which leads to the scalar constraints
\begin{subequations}\label{eq_composition_constraints}
\begin{align}
    G &= -B,\\
    H &= -C,\\
    A &= D,\label{eq_AeqD}\\
    A^{2} - BC &= 1.
\end{align}
\end{subequations}

\subsection{Invariance of the Spacetime Norm}
Define the $(p+q)$-position vector in $\mathrm{O}$ by
\begin{equation}\label{eq_pq_position_def}
    \boldsymbol{\mathcal{X}} \coloneqq \mathbf{x}_{t} + \mathbf{x}_{s},\quad
    \mathbf{x}_{t} = x^{t}\mathbf{e}_{t},\quad
    \mathbf{x}_{s} = x^{s}\mathbf{e}_{s},
\end{equation}
so that applying the direct transformation together with $A = D$ gives
\begin{equation}
    \boldsymbol{\mathcal{X}}
    = \left(Ax'^{t} + Bx'^{s}\right)\mathbf{e}_{t}
    + \left(Cx'^{t} + Ax'^{s}\right)\mathbf{e}_{s},
\end{equation}
whose square produces the real quadratic form
\begin{equation}\label{eq_expansion_direct}
\begin{split}
    \boldsymbol{\mathcal{X}}^{2}
    &= (x^{t})^{2} + \varsigma(x^{s})^{2}= (A^{2}+ \varsigma C^{2})(x'^{t})^{2}\\
    &+ (B^{2}+ \varsigma A^{2})(x'^{s})^{2}
    + 2A(B + \varsigma C)x'^{t}x'^{s}.
\end{split}
\end{equation}
The inverse transformation likewise gives
\begin{equation}\label{eq_expansion_inverse}
\begin{split}
    \boldsymbol{\mathcal{X}}'^{2}
    &= (x'^{t})^{2} + \varsigma(x'^{s})^{2}= (A^{2} + \varsigma C^{2})(x^{t})^{2}\\
    &+ (B^{2} + \varsigma A^{2})(x^{s})^{2}
    - 2A(B + \varsigma C)x^{t}x^{s}.
\end{split}
\end{equation}
Since $\boldsymbol{\mathcal{X}}^{2}$ and $\boldsymbol{\mathcal{X}}'^{2}$ are real scalar norms of the same displacement, reciprocity and isotropy under the Principle of Relativity require
\begin{equation}
    \boldsymbol{\mathcal{X}}^{2} = \boldsymbol{\mathcal{X}}'^{2},
\end{equation}
so equating the coefficients yields
\begin{subequations}\label{eq_scalar_invariance_constraints}
\begin{align}
    A^{2} + \varsigma C^{2} &= 1,\\
    B^{2} + \varsigma A^{2} &= \varsigma,\label{eq_scalar_invariant_for_B}\\
    B + \varsigma C &= 0.
\end{align}
\end{subequations}
For $\varsigma = +1$, we have $B = -C$, and Eq.~\eqref{eq_scalar_invariant_for_B} reduces to
\begin{equation}
    A^{2} + B^{2} = 1,
\end{equation}
which confines the transformation to the circular trigonometric parameterization
\begin{equation}
    A(\varphi) = \cos\varphi,\quad B(\varphi) = \sin\varphi.
\end{equation}
For $\varsigma = -1$, by contrast, $B = C$, and Eq.~\eqref{eq_scalar_invariant_for_B} yields
\begin{equation}
    A^{2} - B^{2} = 1,
\end{equation}
with the hyperbolic parameterization
\begin{equation}
    A(\phi) = \cosh\phi \ge 1,\quad B(\phi) = \sinh\phi.
\end{equation}
Once the spatial and temporal axis directions have been selected, the pure inertial comparisons considered here vary only the relative translational motion, preserving those selections under successive collinear comparisons. However, as the unrestricted parameter $\varphi$ ranges over $\mathbb{R}$, the circular solution carries $A=D$ through positive, zero, and negative values. For $A<0$, successive ticks of a clock at rest in one frame receive decreasing time coordinates in the other, while the signed spatial separation of rod endpoints likewise reverses at the other frame's simultaneity\cite{Data2022_RevisitLorentzTransform}, introducing axis reversals into a comparison intended to describe relative motion alone.

Restricting the circular family to $A>0$ also fails to preserve closure under composition. Indeed, two transformations with $\varphi=\pi/3$ each have $A>0$, whereas their composition has parameter $2\pi/3$ and $A=-1/2$, so two valid changes of inertial frame would together produce an invalid. In contrast, the hyperbolic family preserves $A=D=\cosh\phi\ge1$ throughout $\phi\in\mathbb{R}$ and under arbitrary composition of collinear boosts. Thus, physical admissibility uniquely selects $\varsigma=-1$.

Under this signature, the invariant in a $(p+q)$-dimensional spacetime takes the form
\begin{equation}\label{eq_lorentz_interval}
    \boldsymbol{\mathcal{X}}^{2} = (x^{t})^{2} - (x^{s})^{2},
\end{equation}
which for $p=1$ and $q=3$ is the standard four-dimensional Lorentzian interval $\mathbf{X}^{2}$. In the maximally symmetric case $p=q=3$, the temporal component $\mathbf{x}_{t}$ can rotate while retaining its magnitude $x^{t}$, leaving the interval unchanged. Thus, at fixed $x^{t}$ and $x^{s}$, quantities obtained solely from this quadratic interval are independent of the choice of $p$ and $q$, while that choice determines the directions available to $\boldsymbol{\mathcal{X}}$ in the underlying linear space.

\subsection{The Invariant Limiting Speed}
Tracking the spatial origin of $\mathrm{O}'$ as described in $\mathrm{O}$ identifies the usual scalar relativistic parameters
\begin{equation}
    A = D = \gamma,\quad B = C = \beta\gamma,\quad \beta = \frac{u}{c},
\end{equation}
where $\gamma = (1 - \beta^{2})^{-1/2}$ is the Lorentz factor, so the boost in the selected temporal--spatial plane of a $p$-temporal plus $q$-spatial spacetime becomes
\begin{equation}\label{eq_extended_lorentz_transform}
    \begin{bmatrix}\mathbf{x}_{t}\\\mathbf{x}_{s}\end{bmatrix}
    = \gamma
    \begin{bmatrix}
        \mathbf{e}_{t} & \beta\mathbf{e}_{t}\\
        \beta\mathbf{e}_{s} & \mathbf{e}_{s}
    \end{bmatrix}
    \begin{bmatrix}x'^{t}\\x'^{s}\end{bmatrix}.
\end{equation}
Accordingly, the temporal and spatial vector transformations share the Lorentzian algebraic form, while $\mathbf{e}_{t}$ remains an a priori dynamical direction within the temporal subspace. Because $\beta = u/c$, real and finite boosts require $|u| < c$, identifying $c$ with the invariant limiting speed and, experimentally, with the speed of light in vacuum.

Because the spacetime norm is invariant, and the speed of light is universal in any $(p+q)$-dimensional spacetime obeying relativity, the usual definition of proper time,
\begin{equation}
\label{eq_pqvector_squared}
    \boldsymbol{\mathcal{X}}^{2}=\boldsymbol{\mathcal{X}} \cdot \boldsymbol{\mathcal{X}} = c^{2}\tau^{2},
\end{equation}
extends to these spaces, allowing for a generalized $(p+q)$-dimensional velocity defined as
\begin{equation}
    \boldsymbol{\mathcal{U}} \coloneqq \frac{\mathrm{d}\boldsymbol{\mathcal{X}}}{\mathrm{d}\tau},
\end{equation}
with norm $|\boldsymbol{\mathcal{U}}| = c$.

To establish this result, we first define a vector $\boldsymbol{\mathcal{V}}$ norm as
\begin{equation}
    |\boldsymbol{\mathcal{V}}| \coloneqq \sqrt{\boldsymbol{\mathcal{V}}^{2}},
\end{equation}
and differentiate $|\boldsymbol{\mathcal{X}}|^{2} = \boldsymbol{\mathcal{X}}^{2}$ for the position vector $\boldsymbol{\mathcal{X}}$ with respect to $\tau$ using the product rule on the inner product $\boldsymbol{\mathcal{X}}^{2} = \boldsymbol{\mathcal{X}} \cdot \boldsymbol{\mathcal{X}}$, which yields the kinematic relation
\begin{equation}
    \frac{\mathrm{d}}{\mathrm{d}\tau}|\boldsymbol{\mathcal{X}}| = \frac{\boldsymbol{\mathcal{X}} \cdot \boldsymbol{\mathcal{U}}}{|\boldsymbol{\mathcal{X}}|}.
\end{equation}
Substituting $|\boldsymbol{\mathcal{X}}| = c\tau$ back into this expression yields $\boldsymbol{\mathcal{X}} \cdot \boldsymbol{\mathcal{U}} = c |\boldsymbol{\mathcal{X}}| = c^{2}\tau$, which, for a homogeneous inertial trajectory that passes through the origin, gives $\boldsymbol{\mathcal{X}} = \tau \boldsymbol{\mathcal{U}}$, yielding
\begin{equation}\label{eq_squared_pqvelocity_norm}
    \boldsymbol{\mathcal{U}}^{2} = c^{2} \implies |\boldsymbol{\mathcal{U}}| = c.
\end{equation}
Because inertial frames are related by uniform motion, their reference trajectories are straight lines, so along a homogeneous inertial trajectory the finite interval relation in Eq.~\eqref{eq_pqvector_squared} gives the local differential form $(\mathrm{d}\boldsymbol{\mathcal{X}})^{2}=\mathrm{d}\boldsymbol{\mathcal{X}}\cdot\mathrm{d}\boldsymbol{\mathcal{X}}=c^{2}(\mathrm{d}\tau)^{2}$. Moreover, an accelerating observer in flat spacetime has an instantaneously comoving inertial frame at each event, so the local covariant relations established in inertial frames extend pointwise to smoothly accelerated observers. This extension, termed here the \emph{comoving inertial frames corollary}, will be used repeatedly throughout this paper, particularly in the derivation of the compatibility condition and the Planck--de Broglie relations.
